Безразмерное уравнение Рейнольдса~\eqref{eq:reynolds_nondim}
в зоне полного заполнения смазочного слоя линейно по давлению:
коэффициент~$\bar{h}^{\,3}$ зависит только от геометрии
и не содержит неизвестных. При наличии микроструктур
поверхности линейность по $p$ сохраняется, однако
вычислительная сложность задачи существенно возрастает~\cite{miraskari2018a,miraskari2018b}.
Во-первых, локальные перепады зазора в ячейках микрорельефа
создают резкий контраст коэффициентов: кубическая
зависимость~$h^3$ усиливает разницу между глубокими и мелкими
участками зазора, что ухудшает обусловленность дискретной
системы и замедляет сходимость итерационных решателей
(раздел~3.3). Во-вторых, кавитация в зонах расходящегося
зазора порождает негладкие ограничения~-- переключение между
зонами полного заполнения и кавитации, совместное определение
давления~$p$ и степени заполнения~$\vartheta$ через условия
комплементарности~\eqref{eq:jfo_complementarity}. В-третьих,
при определённых условиях плотность смазки зависит от искомого
давления ($\rho = \rho(p)$), что вносит дополнительную, но уже
гладкую нелинейность. Ниже каждый из перечисленных эффектов
рассматривается отдельно.

\textit{Кавитация как источник негладкой нелинейности.}
Физическая природа кавитации и массосохраняющая постановка JFO
подробно описаны в подразделах~2.3.1-2.3.2. Здесь существенно
то, что постановка JFO~-- это задача с условиями
комплементарности~\eqref{eq:jfo_complementarity}: в каждом узле
либо $p > p_{\mathrm{cav}}$ и $\vartheta = 1$, либо
$p = p_{\mathrm{cav}}$ и $\vartheta < 1$, причём граница между
зонами заранее неизвестна. Вместо единой линейной системы
возникает задача с неравенствами, которая решается итеративно
методом активного множества~-- внешний цикл обновляет
классификацию узлов и пересчитывает~$\vartheta$, а внутренний
цикл решает линейную подсистему для давления при
фиксированном~$\vartheta$ (глава~3, раздел~3.5). Такая
двухуровневая итерационная структура требует согласования
критериев остановки обоих уровней.

При гладком зазоре кавитация, как правило, формирует одну
связную область, положение которой стабилизируется за несколько
внешних итераций. Микроструктуры поверхности принципиально
усложняют картину: каждый элемент микрорельефа создаёт
локальный перепад зазора с зоной разрежения на расходящемся
участке. В результате кавитационная область может формироваться
не в виде одной связной зоны, а в виде множества
фрагментированных «островков», привязанных к отдельным ячейкам
микрорельефа. Активное множество узлов
с~$p = p_{\mathrm{cav}}$ становится разрозненным, его
конфигурация чувствительна к разрешению сетки и к критериям
остановки внешних итераций. Ошибка в определении границы
кавитации даже на один узел изменяет баланс потоков в соседних
ячейках, поэтому требуется достаточное разрешение сетки для
описания каждого элемента микрорельефа (не менее нескольких
узлов на ячейку) и ужесточение критериев сходимости внешнего
цикла по сравнению с задачей для гладкого зазора.

\textit{Роль параметра $\Lambda$ и анизотропия задачи.}
Параметр
$\Lambda = L_1/L_2$~\eqref{eq:nondim_Lambda} определяет
соотношение вкладов продольного ($s_1$) и поперечного ($s_2$)
членов в уравнении~\eqref{eq:reynolds_nondim}:
при $\Lambda \ll 1$ вклад второго члена мал и задача
приближается к одномерной по~$s_1$, при $\Lambda \gtrsim 1$
существенно возрастает роль поперечного растекания по~$s_2$
(коэффициент перед соответствующим членом
равен $\Lambda^2$). Для задач с микроструктурами
значение~$\Lambda$ определяет, насколько локальные эффекты
микрорельефа «размываются» потоком в поперечном направлении.
При большом~$\Lambda$ поперечное растекание сглаживает
возмущения давления от отдельных ячеек; при малом~$\Lambda$
эти возмущения концентрируются вдоль~$s_1$ и оказывают более
сильное влияние на результирующее поле давления и границы
кавитации.

Конкретные значения~$\Lambda$ для трёх постановок монографии
следуют из масштабов таблицы~\ref{tab:nondim_scales}.
Для радиального подшипника $L_1 = R$ (окружной масштаб),
$L_2 = L$ (длина подшипника), и значения $\Lambda = R/L$
соответствуют умеренной или сильной анизотропии в зависимости
от геометрии. Для упорного подшипника $L_1 = R$,
$L_2 = R_{\mathrm{out}} - R_{\mathrm{in}}$
и $\Lambda = R/(R_{\mathrm{out}} - R_{\mathrm{in}})$, как
правило, порядка единицы, т.\,е.\ задача квази-изотропна.
Для возвратно-поступательного подшипника $L_1 = L$ (длина
цилиндра), $L_2 = R$ и $\Lambda = L/R$ может быть значительно
больше единицы; это означает усиление вклада поперечного члена
(по $s_2 = R\theta$) в безразмерном
уравнении~\eqref{eq:reynolds_nondim} и повышенные требования
к разрешению по поперечному направлению.

С вычислительной точки зрения анизотропия задачи требует
согласования шагов сетки по~$s_1$ и~$s_2$. При больших
значениях $\Lambda$ вклад производных по~$s_2$ усиливается
коэффициентом~$\Lambda^2$, и недостаточное разрешение
по «слабому» направлению может приводить
к дискретизационным ошибкам, доминирующим в оценке сходимости
и искажающим вычисленные границы кавитационных зон. Разумная
стратегия~-- выбирать отношение
$\Delta s_1/\Delta s_2$ порядка~$\Lambda$, чтобы
безразмерные шаги сетки были сопоставимы.

\textit{Изменение плотности: гладкая нелинейность
$\rho = \rho(p)$.}
В задачах глав~5-7 смазка предполагается несжимаемой
жидкостью, $\rho = \mathrm{const}$, и плотность не входит
в число неизвестных. Вместе с тем существуют ситуации, когда
это допущение неприемлемо: газовая смазка, большие перепады
давления в высоконагруженных узлах, двухфазные режимы
со значительным изменением массового содержания фаз в области
кавитации. Для полноты постановки здесь приводится обобщение
уравнения JFO~\eqref{eq:jfo_conservation} на случай переменной
плотности в размерной форме:
\begin{equation}
  \frac{\partial (\rho\,\vartheta\,h)}{\partial t}
  + \nabla\cdot\!\left(
      -\frac{\rho\,h^3}{12\mu}\,\nabla p
      + \frac{\rho\,\vartheta\,h}{2}\,\mathbf{U}
    \right) = 0,
  \label{eq:jfo_rho_general}
\end{equation}
где плотность~$\rho = \rho(p)$ входит как в член накопления,
так и в напорный и переносной потоки. В отличие от
комплементарности JFO, зависимость~$\rho(p)$ является гладкой
(задаётся уравнением состояния конкретной среды), поэтому результирующая нелинейность также гладкая
и устраняется итерациями Пикара с релаксацией: на каждой
итерации $\rho$ пересчитывается по полю давления предыдущего
шага, после чего решается обновлённая система. Такой подход
совместим с двухуровневой структурой алгоритма JFO
(раздел~3.5) и не требует изменения внутреннего
решателя~-- достаточно обновлять коэффициенты системы перед
очередной внешней итерацией. Следует подчеркнуть, что
в расчётах глав~5-7 данной монографии
$\rho = \mathrm{const}$
и обобщение~\eqref{eq:jfo_rho_general} не используется;
раздел описывает его для полноты постановки.

\textit{Термовязкость и внешние итерации.}
Гладкая нелинейность возникает и при учёте термовязкостных
эффектов: если вязкость зависит от температуры
($\mu = \mu(T)$), а температура в свою очередь определяется
скоростью диссипации в зазоре, то уравнение Рейнольдса
оказывается связано с уравнением энергии. В этом случае задача
теряет линейность по $p$: при фиксированном поле $T$ (а значит,
$\mu = \mu(T)$) уравнение по-прежнему линейно по давлению,
однако $T$ зависит от градиентов давления и скорости, что
требует итерационного согласования. Типичная стратегия~--
внешний цикл Пикара (fixed-point): на каждой итерации сначала
решается уравнение Рейнольдса при фиксированном поле вязкости,
затем из полученного поля скоростей пересчитывается температура
и вязкость, и цикл повторяется до выхода на согласованное
решение. Сходимость внешнего цикла, как правило, линейная;
число итераций составляет 5--20 в зависимости от интенсивности
тепловых эффектов. В изотермической постановке, принятой
в прикладных главах настоящей монографии, внешний цикл
не требуется: уравнение Рейнольдса решается как линейная
задача при фиксированной $\mu = \mu_0 = \mathrm{const}$.

Подводя итог, наиболее принципиальная нелинейность при
моделировании смазочного слоя с микроструктурами~-- кавитация
в постановке JFO, порождающая задачу с условиями
комплементарности и фрагментированным активным множеством.
Параметр~$\Lambda$ характеризует анизотропию уравнения
Рейнольдса и режим растекания, определяя степень влияния
локальных эффектов микрорельефа на общую картину давления
и конфигурацию кавитационных зон. Зависимость $\rho(p)$
представляет собой опциональное обобщение, применяемое при
наличии физических предпосылок; в расчётах глав~5-7
монографии смазка принимается несжимаемой~\cite{jang2002,rebufa2017}.
